% Frontmatter: abstract, contributions, reproducibility

\begin{abstract}
给定有限窗口本体空间 $\Omega_m=\{0,1\}^m$ 与稳定输出语言 $X_m\subset\{0,1\}^m$（禁止子词 $11$），
折叠映射 $\Fold_m:\Omega_m\twoheadrightarrow X_m$ 将窗口微态规范化为 Zeckendorf 约束下的稳定词。
本文聚焦于一个更强的问题：当我们不仅需要输出 $\Fold_m(b)$，还希望保留最少的附加信息使得后续任务
（例如可逆化、搜索或预测接口）可被有限状态实现时，\emph{最小残差}应如何定义并如何计算？

本文提出一个以 Myhill--Nerode 右等价为核心的残差最小化框架：把“对未来输入的响应”视为右作用下的等价类，
从而将残差集合刻画为最小有限状态结构的状态集。作为主结果之一，我们对本文所定义的 $\Fold_m$ 给出残差规模谱
闭式与显式整数状态同构；进一步给出一个一般化模板：在“截断编码在某个整数区间上保持单射”的条件下，
深度 $k$ 的右等价类与前缀加权和一一对应，从而把残差计数归约为子集和计数问题。
\end{abstract}

\paragraph{贡献概览}
\begin{itemize}
  \item (C1) 给出“残差=右等价类”的形式定义，并将其与最小有限状态实现联系起来。
  \item (C2) 给出残差规模谱的闭式结果，并给出由 Fibonacci 前缀和诱导的显式状态同构与转移更新。
  \item (C3) 给出一般化残差谱模板：在单射窗口条件下，右等价类由前缀加权和刻画，并可归约为子集和计数。
  \item (C4) 提供可复现的枚举与核对协议，输出 $m$ 扫描下的残差规模、纤维谱与熵量等表格产物。
\end{itemize}

\paragraph{可复现性声明}
本稿的数值产物由 \verb|scripts/run_all.py| 生成并写入 \verb|sections/generated/|。
该目录中的 \verb|.tex| 与 \verb|.json| 文件会在编译时被正文自动包含。

